\section{Resultados, Conclusões e Trabalhos Futuros}
\label{sec:resultados}

Esta seção sintetiza o que o trabalho produziu, delimita o alcance das conclusões e indica como ele pode ser continuado.

\subsection{Resultados experimentais}

\pendente{Esta subseção depende da execução dos blocos A--D e do micro-experimento. Substituir os marcadores pelos valores obtidos, inserir as figuras geradas por \texttt{analysis.py} (curvas de convergência por $\alpha$; mapa de calor $\alpha \times \mu$) e reescrever os parágrafos abaixo à luz dos dados.}

A Tabela~\ref{tab:resultados} apresenta, para o CIFAR-10, a acurácia final de FedAvg e FedProx com $\mu^*$ nas dez sementes confirmatórias, a diferença pareada com intervalo de confiança, o p-valor do Wilcoxon ajustado por Holm, o p-valor do teste de equivalência (TOST) e as rodadas até a acurácia-alvo. A mesma tabela para o Fashion-MNIST e os resultados dos algoritmos de referência ficam no repositório.

\begin{table}[ht]
    \centering\scriptsize
    \caption{FedAvg × FedProx ($\mu^* = $ \pendente{?}) no CIFAR-10, bloco D (10 sementes). Acurácia em \%, média $\pm$ desvio padrão.}
    \label{tab:resultados}
    \setlength{\tabcolsep}{3pt}
    \begin{tabular}{@{}lcccccc@{}}
        \toprule
        $\alpha$ & FedAvg & FedProx & $\Delta$ [IC 95\%] & $p_{\text{Holm}}$ & $p_{\text{TOST}}$ & $R_\tau$ (Avg/Prox) \\
        \midrule
        0,05 & \pendente{--} & \pendente{--} & \pendente{--} & \pendente{--} & \pendente{--} & \pendente{--} \\
        0,1  & \pendente{--} & \pendente{--} & \pendente{--} & \pendente{--} & \pendente{--} & \pendente{--} \\
        0,5  & \pendente{--} & \pendente{--} & \pendente{--} & \pendente{--} & \pendente{--} & \pendente{--} \\
        1    & \pendente{--} & \pendente{--} & \pendente{--} & \pendente{--} & \pendente{--} & \pendente{--} \\
        100  & \pendente{--} & \pendente{--} & \pendente{--} & \pendente{--} & \pendente{--} & \pendente{--} \\
        \bottomrule
    \end{tabular}
\end{table}

\textbf{QP1 (H1--H3).} \pendente{Descrever em que faixa de $\alpha$ o ganho do FedProx é significativo, se o melhor $\mu$ muda com $\alpha$ e se o ranking muda entre partição patológica e desvio de quantidade. Comparar com FedAvgM, FedAdam e SCAFFOLD.}

\textbf{QP2 (H4).} \pendente{Informar a razão entre as medianas do tempo por época (FedProx/FedAvg) na CPU de 1 thread, a diferença de memória de pico e o custo total até o alvo $C_\tau$ para cada $\alpha$.}

\subsection{Principais contribuições}

Em relação aos objetivos da Seção~\ref{sec:introducao}, o trabalho entrega:

\begin{itemize}[leftmargin=*, itemsep=2pt]
    \item \textbf{Levantamento sistemático focado em estudos comparativos (OE1).} São 22 estudos de 2021 a 2026, selecionados por critérios explícitos e reprodutíveis, com protocolo, tabelas por critério e planilha de triagem disponíveis no repositório. A análise mostrou que a literatura se contradiz sobre o FedProx e identificou três lacunas: $\mu$ não controlado, custo local não medido e ausência de um limiar de heterogeneidade.
    \item \textbf{Ambiente experimental reprodutível e de custo zero (OE2, OE3).} Inclui simulação com Flower e partições controladas, pareamento por semente e teste de sanidade $\mu = 0 \equiv$ FedAvg, tudo executável em GPUs gratuitas. O ambiente pode ser reaproveitado para avaliar outros algoritmos de agregação sem mudança no restante do \textit{pipeline}.
    \item \textbf{Evidência sobre quando o FedProx compensa (OE4--OE6).} \pendente{Resumir o mapa $\alpha \times \mu \times$ (acurácia, custo) obtido e a recomendação prática derivada, por exemplo: ``para $\alpha \leq$ X, usar FedProx com $\mu =$ Y; acima disso, FedAvg ou FedAvgM''.}
    \item \textbf{Medição isolada do custo do termo proximal em CPU restrita.} \pendente{Resumir o resultado do micro-experimento.} Nenhum dos estudos do levantamento relata essa medição no resumo.
\end{itemize}

\subsection{Limitações do trabalho}

\begin{itemize}[leftmargin=*, itemsep=2pt]
    \item \textbf{Simulação × implantação.} Todos os clientes rodam em uma única máquina, sem latência de rede, falhas de conexão nem diferenças reais de hardware. \citet{valli2026beyond} mostram que resultados de simulação podem divergir de implantações reais. O micro-experimento em CPU de um núcleo aproxima o custo de um dispositivo de borda, mas não o substitui.
    \item \textbf{Modelos e dados.} As CNNs são pequenas e os datasets são de imagens com dez classes balanceadas. Redes profundas, outros domínios (séries temporais, texto) ou desvio de atributos podem se comportar de outra forma. As conclusões valem para o regime estudado.
    \item \textbf{Hiperparâmetros fixos.} A taxa de aprendizado local e o número de épocas são os mesmos para todos os algoritmos e não foram otimizados para cada um. É uma escolha deliberada para isolar o efeito da agregação, mas pode favorecer algum algoritmo.
    \item \textbf{Grade de $\mu$ e orçamento.} Foram testados três valores de $\mu$, e os algoritmos de referência rodaram com apenas três sementes, por causa da cota gratuita de GPU.
    \item \textbf{Cobertura do levantamento.} Scopus e Web of Science não foram consultadas diretamente, e a triagem foi feita por título e resumo. Quatro estudos com resumo inacessível ficaram fora da análise.
\end{itemize}

\subsection{Possibilidades de continuidade}

As extensões mais diretas atacam as limitações acima:

\begin{itemize}[leftmargin=*, itemsep=2pt]
    \item \textbf{Validação em hardware real.} Executar FedAvg e FedProx com o mesmo código em um conjunto de Raspberry Pi ou celulares, usando o modo de implantação do Flower, e medir tempo, energia e tráfego de rede reais, como em \citet{carranco2025adhoc}.
    \item \textbf{Heterogeneidade de sistema.} Executar o cenário C5 (\textit{stragglers}), em que o FedProx aceita trabalho local parcial e onde, segundo \citet{cheng2026stragglers}, as diferenças entre os algoritmos são maiores.
    \item \textbf{$\mu$ adaptativo.} Se o melhor $\mu$ variar com $\alpha$, investigar regras que ajustem $\mu$ por rodada ou por cliente a partir de uma estimativa local da heterogeneidade.
    \item \textbf{Desvio de atributos e dados naturais.} Incluir o FEMNIST \citep{caldas2018leaf}, com heterogeneidade natural por escritor, e datasets de domínios de borda como sensores IoT e reconhecimento de atividades.
    \item \textbf{Mais algoritmos e modelos.} Estender a comparação a FedNova \citep{wang2020fednova} e a modelos maiores (ResNet), verificando se o custo relativo do termo proximal muda com o tamanho da rede.
\end{itemize}

\subsection{Aplicações futuras}

O resultado principal, um critério de escolha do algoritmo de agregação a partir do grau de heterogeneidade, é útil a quem projeta sistemas federados em borda. Alguns exemplos:

\begin{itemize}[leftmargin=*, itemsep=2pt]
    \item \textbf{Saúde conectada:} hospitais e dispositivos vestíveis com populações de pacientes muito diferentes, domínio mais frequente no levantamento \citep{mahmood2025alzheimer, mim2025healthcare}.
    \item \textbf{IoT industrial e cidades inteligentes:} sensores e câmeras fixos que observam fenômenos locais, como no reconhecimento de vagas \citep{bakr2026parkfl} e na previsão de geração renovável \citep{ghosh2026renewable}.
    \item \textbf{Enxames de drones e veículos:} clientes móveis com dados que mudam de distribuição conforme a missão \citep{hu2025uav}.
\end{itemize}

Em todos esses cenários, uma estimativa simples da heterogeneidade antes da implantação (por exemplo, a divergência entre as distribuições de rótulos dos clientes) pode ser cruzada com o mapa produzido aqui para decidir se o custo extra do FedProx nos dispositivos se justifica. O repositório aberto permite que essa avaliação seja refeita para o modelo e os dados de cada aplicação.

\subsection{Conclusão}

O trabalho partiu de uma constatação do levantamento bibliográfico: apesar de amplamente usado, o FedProx não tem uma vantagem consistente sobre o FedAvg na literatura recente, e os estudos raramente controlam $\mu$, o nível de heterogeneidade ou o custo no cliente. Para responder em quais condições o termo proximal compensa, foi desenhado um experimento pareado, reprodutível e de custo zero, que cruza tipo e nível de heterogeneidade com o valor de $\mu$ e mede tanto a qualidade do modelo quanto o custo local. \pendente{Concluir com a resposta às QP1 e QP2 em uma ou duas frases, com base nos resultados.}
